\documentclass[11pt]{article}
\usepackage[utf8]{inputenc}
\usepackage{amsmath,amssymb}
\usepackage{hyperref}
\title{Robust Adversarial Malware Detection under Distribution Shift}
\author{Anonymous}
\date{}
\begin{document}
\maketitle

\begin{abstract}
This paper studies Adversarial Malware Detection. Grounded in a real, citable literature \cite{ref1,ref2,ref3,ref4,ref5,ref6,ref7,ref8},
we propose a focused method and report \textbf{illustrative} experiments.
All references below are real and verifiable (OpenAlex/arXiv); the reported
numbers are simulated to illustrate intended behaviour.
\end{abstract}

\section{Introduction}
Adversarial Malware Detection is an active research area. This work targets a concrete gap identified from
the recent literature and contributes a method together with an evaluation protocol.

\section{Related Work (real literature)}
The following works, retrieved from public bibliographic APIs, frame our study:
\begin{itemize}
\item Biggio et al. (2013) — "Evasion attacks against machine learning at test time" (cited 1636).
\item Grosse et al. (2017) — "Adversarial Examples for Malware Detection", Lecture notes in computer science (cited 555).
\item Luis et al. (2017) — "Towards poisoning of deep learning algorithms with back-gradient optimization", UNICA IRIS Institutional Research Information System (University of Cagliari) (cited 541).
\item Grosse et al. (2017) — "On the (Statistical) Detection of Adversarial Examples", arXiv (Cornell University) (cited 377).
\item Papernot et al. (2018) — "Deep k-Nearest Neighbors: Towards Confident, Interpretable and Robust Deep Learning", arXiv (Cornell University) (cited 370).
\item Hu et al. (2017) — "Generating Adversarial Malware Examples for Black-Box Attacks Based on GAN", arXiv (Cornell University) (cited 366).
\end{itemize}

\section{Method}
We model distribution shift explicitly and optimise a worst-case objective over an uncertainty set of plausible test distributions. We instantiate this idea for Adversarial Malware Detection and analyse its cost/quality trade-off.

\section{Experiments (Illustrative)}
\begin{center}
\begin{tabular}{lcc}
System & Primary Metric & Efficiency \\ \hline
Strong baseline & 0.812 & 1.00$\times$ \\
Ablation & 0.828 & 0.92$\times$ \\
\textbf{Ours} & \textbf{0.851} & \textbf{0.71$\times$}
\end{tabular}
\end{center}
\emph{Metrics simulated to illustrate the intended effect; not measured.}

\section{Conclusion}
We connected a real literature to a concrete method for Adversarial Malware Detection. Future work will
validate the illustrative results with real experiments.

\begin{thebibliography}{9}
\bibitem{ref1} Battista Biggio, Igino Corona, Davide Maiorca et al.. Evasion attacks against machine learning at test time. \emph{}, 2013. 
\bibitem{ref2} Kathrin Grosse, Nicolas Papernot, Praveen Manoharan et al.. Adversarial Examples for Malware Detection. \emph{Lecture notes in computer science}, 2017. DOI:10.1007/978-3-319-66399-9\_4
\bibitem{ref3} Muñoz-González, Luis, Battista Biggio, Ambra Demontis et al.. Towards poisoning of deep learning algorithms with back-gradient optimization. \emph{UNICA IRIS Institutional Research Information System (University of Cagliari)}, 2017. DOI:10.1145/3128572.3140451
\bibitem{ref4} Kathrin Grosse, Praveen Manoharan, Nicolas Papernot et al.. On the (Statistical) Detection of Adversarial Examples. \emph{arXiv (Cornell University)}, 2017. DOI:10.48550/arxiv.1702.06280
\bibitem{ref5} Nicolas Papernot, Patrick McDaniel. Deep k-Nearest Neighbors: Towards Confident, Interpretable and Robust Deep Learning. \emph{arXiv (Cornell University)}, 2018. DOI:10.48550/arxiv.1803.04765
\bibitem{ref6} Weiwei Hu, Ying Tan. Generating Adversarial Malware Examples for Black-Box Attacks Based on GAN. \emph{arXiv (Cornell University)}, 2017. DOI:10.48550/arxiv.1702.05983
\bibitem{ref7} Xiao Chen, Chaoran Li, Derui Wang et al.. Android HIV: A Study of Repackaging Malware for Evading Machine-Learning Detection. \emph{arXiv (Cornell University)}, 2018. DOI:10.1109/tifs.2019.2932228
\bibitem{ref8} Kathrin Grosse, Nicolas Papernot, Praveen Manoharan et al.. Adversarial Perturbations Against Deep Neural Networks for Malware Classification. \emph{arXiv (Cornell University)}, 2016. DOI:10.48550/arxiv.1606.04435
\end{thebibliography}
\end{document}
